% !TEX root = ../main.tex
\chapter{Higher-Order Differentiation}
\label{ch:higher_order}

The basic autograd machinery computes first-order gradients: the
sensitivity of a scalar output to its inputs. Many algorithms ---
second-order optimisers, sensitivity analyses, certain
regularisers, some interpretability methods --- need higher-order
information: gradients of gradients, Jacobians, Hessians,
Hessian-vector products. \Lucid\ exposes a small set of
\emph{functional transforms} that compose on top of the autograd
engine to deliver these.

The transforms (\code{grad}, \code{vjp}, \code{jvp}, \code{vmap},
\code{jacrev}, \code{jacfwd}, \code{hessian}) are imported from
JAX-style functional differentiation. They are exposed under
\code{lucid.func} and live in \code{lucid/func/}. This chapter
describes each, how they compose, and the implementation in terms
of the underlying autograd engine.

\section{Overview}
\label{sec:higher_order:overview}

The functional transforms take a Python function and return a
new Python function:

\begin{itemize}
  \item \code{grad(f)}: returns a function that computes the
    gradient of \code{f}.
  \item \code{vjp(f, *primals)}: returns the value and a
    vjp-function that applies $J^\top$.
  \item \code{jvp(f, primals, tangents)}: returns the value and
    the jvp ($J$ applied to tangents).
  \item \code{vmap(f)}: returns a function that maps over the
    leading axis.
  \item \code{jacrev(f)}: returns a function that computes the
    Jacobian via reverse mode.
  \item \code{jacfwd(f)}: returns a function that computes the
    Jacobian via forward mode.
  \item \code{hessian(f)}: returns a function that computes the
    Hessian (typically \code{jacrev(jacrev(f))} or
    \code{jacfwd(jacrev(f))}).
\end{itemize}

All compose. \code{grad(vmap(f))} computes per-sample gradients
in parallel. \code{hessian = jacfwd(jacrev(f))} computes the
Hessian via mixed mode.

\section{grad}
\label{sec:higher_order:grad}

\code{grad(f)} returns a function that, given the same arguments
as \code{f}, returns the gradient of \code{f}'s scalar output
with respect to the first argument.

\begin{lstlisting}[style=lucid,language=LucidPython]
import lucid
from lucid.func import grad

def loss(w, x, y):
    pred = w @ x
    return ((pred - y) ** 2).sum()

grad_loss = grad(loss)
w = lucid.randn(3, 5)
x = lucid.randn(5)
y = lucid.randn(3)
g = grad_loss(w, x, y)   # gradient of loss with respect to w
\end{lstlisting}

The output \code{g} is a tensor of the same shape as \code{w},
holding $\nabla_w \text{loss}(w, x, y)$.

\subsection{grad over a different argument}
\label{ssec:higher_order:grad_argnum}

By default \code{grad} differentiates with respect to the first
argument. The \code{argnums} parameter selects which argument:

\begin{lstlisting}[style=lucid,language=LucidPython]
grad_x = grad(loss, argnums=1)   # gradient w.r.t. x (the input)
\end{lstlisting}

A tuple of \code{argnums} returns a tuple of gradients:

\begin{lstlisting}[style=lucid,language=LucidPython]
grad_wx = grad(loss, argnums=(0, 1))
gw, gx = grad_wx(w, x, y)
\end{lstlisting}

\subsection{Implementation}
\label{ssec:higher_order:grad_impl}

\code{grad(f)} returns a wrapper that:

\begin{enumerate}
  \item Detaches the inputs and marks the differentiated argument
    as \code{requires\_grad=True}.
  \item Calls \code{f} with the prepared inputs, getting a
    scalar output.
  \item Calls \code{output.backward(create\_graph=False)}.
  \item Returns the differentiated argument's \code{.grad}.
\end{enumerate}

The wrapping is around 30 lines of Python; the heavy lifting is
the same autograd engine that \code{tensor.backward()} uses.

\section{vjp and jvp}
\label{sec:higher_order:vjp_jvp}

These are the low-level building blocks for the other transforms.

\subsection{vjp}
\label{ssec:higher_order:vjp}

\code{vjp(f, *primals)} returns a tuple \code{(output,
vjp\_fn)}. \code{output} is \code{f(*primals)}. \code{vjp\_fn} is
a function that, given a cotangent of the output, returns the
cotangent of the inputs:

\begin{lstlisting}[style=lucid,language=LucidPython]
from lucid.func import vjp

def f(x):
    return x.sin() + x.cos()

y, vjp_fn = vjp(f, x)
g, = vjp_fn(lucid.ones_like(y))    # gradient of sum-of-output w.r.t. x
\end{lstlisting}

The pattern is useful when the same forward computation will be
differentiated multiple times with different upstream cotangents
(e.g., for computing different rows of a Jacobian).

\subsection{jvp}
\label{ssec:higher_order:jvp}

\code{jvp(f, primals, tangents)} returns \code{(output,
jvp\_value)}. \code{jvp\_value} is $J \cdot \text{tangents}$, the
directional derivative of \code{f} at \code{primals} in the
direction of \code{tangents}:

\begin{lstlisting}[style=lucid,language=LucidPython]
from lucid.func import jvp

y, dy = jvp(f, (x,), (lucid.ones_like(x),))
\end{lstlisting}

\code{jvp} is forward-mode AD; it does not require building a
backward graph.

\section{vmap}
\label{sec:higher_order:vmap}

\code{vmap(f)} returns a function that applies \code{f} to a
batch of inputs in parallel, treating the leading axis of each
input as the batch dimension.

\begin{lstlisting}[style=lucid,language=LucidPython]
from lucid.func import vmap

def per_sample_loss(w, x, y):
    return ((w @ x - y) ** 2).sum()

# Apply per-sample loss across a batch:
batched_loss = vmap(per_sample_loss, in_axes=(None, 0, 0))
total = batched_loss(w, x_batch, y_batch).sum()
\end{lstlisting}

The \code{in\_axes} argument names which axis of each input to
vectorise over. \code{None} means broadcast (no batching for that
argument). \code{0} means batch over the leading axis.

\subsection{Implementation}
\label{ssec:higher_order:vmap_impl}

\code{vmap} is implemented on top of \MLX's own \code{vmap}, which
performs the broadcasting transformation at the primitive level
(\chref{ch:mlx_runtime}). \Lucid's wrapping handles the
\code{Tensor} $\leftrightarrow$ \code{mlx::core::array} unwrap and
the autograd metadata threading.

The composition \code{grad(vmap(f))} is the standard pattern for
per-sample gradients --- vmap over the batch, grad over the
parameter. Equivalent to a manual Python loop but parallelised
inside the kernel.

\section{jacrev and jacfwd}
\label{sec:higher_order:jacrev_jacfwd}

These compute Jacobians via reverse and forward mode
respectively.

\subsection{jacrev}
\label{ssec:higher_order:jacrev}

\code{jacrev(f)} returns a function that computes the full
Jacobian of \code{f} at its argument via reverse mode:

\begin{lstlisting}[style=lucid,language=LucidPython]
from lucid.func import jacrev

def f(x):
    return lucid.stack([x[0]**2, x[1]**3, (x[0] + x[1]).sin()])

J = jacrev(f)(x)    # shape (3, 2): 3 outputs, 2 inputs
\end{lstlisting}

Internally \code{jacrev} runs \code{m} VJP calls (where \code{m}
is the output dimension), each with a one-hot cotangent. The
results are stacked into the Jacobian.

\subsection{jacfwd}
\label{ssec:higher_order:jacfwd}

\code{jacfwd(f)} computes the same Jacobian via forward mode,
running \code{n} JVP calls (where \code{n} is the input dimension):

\begin{lstlisting}[style=lucid,language=LucidPython]
from lucid.func import jacfwd

J = jacfwd(f)(x)    # same shape, computed via forward mode
\end{lstlisting}

The choice between \code{jacrev} and \code{jacfwd} is the
output-vs-input-dimension trade-off: reverse mode wins when
output is smaller, forward when input is smaller. For square
Jacobians the choice depends on the kernel implementation's
constant factor; \code{jacrev} is usually slightly faster in
\Lucid\ because its building block (VJP) is what the framework
optimises for.

\section{hessian}
\label{sec:higher_order:hessian}

The Hessian of a scalar-output function is the matrix of
second-order partial derivatives. \code{hessian(f)} returns a
function that computes it.

\subsection{Composition}
\label{ssec:higher_order:hessian_composition}

The standard implementation:

\begin{lstlisting}[style=lucid,language=LucidPython]
from lucid.func import jacrev, jacfwd

def hessian(f):
    return jacfwd(jacrev(f))
\end{lstlisting}

The inner \code{jacrev} gives the gradient (a vector); the outer
\code{jacfwd} differentiates each component, producing the
Hessian.

The choice of \code{jacfwd o jacrev} (mixed mode) over
\code{jacrev o jacrev} (pure reverse) is a memory-vs-time
trade-off: pure reverse builds two backward graphs simultaneously
(memory-heavy); mixed mode builds only one (the inner reverse)
plus a forward-mode pass over its output (memory-lighter).

\subsection{Hessian-vector products}
\label{ssec:higher_order:hvp}

For most second-order applications, the full Hessian is too
expensive to compute and store ($O(N^2)$ for $N$ parameters). The
quantity actually needed is the Hessian-vector product
$H v$, which can be computed as the gradient of $v^\top \nabla L$
at cost comparable to one extra backward pass.

\Lucid\ does not expose an \code{hvp} primitive (it is short to
write by hand), but the pattern:

\begin{lstlisting}[style=lucid,language=LucidPython]
def hvp(f, x, v):
    # H @ v via the trick H @ v = grad(v . grad(f))(x)
    g = grad(f)(x)
    return grad(lambda x_: (grad(f)(x_) * v).sum())(x)
\end{lstlisting}

This is the standard mechanism for Newton-style optimisers (which
need $H^{-1} g$, approximated via conjugate gradient on $H v$
products).

\section{Composition Rules}
\label{sec:higher_order:composition}

The transforms compose freely. The composition is the source of
their power.

\subsection{grad o grad}
\label{ssec:higher_order:grad_grad}

The second derivative of a scalar function with respect to a
scalar input:

\begin{lstlisting}[style=lucid,language=LucidPython]
def f(x):
    return x.sin() * x.cos()

f_prime = grad(f)
f_double_prime = grad(f_prime)
\end{lstlisting}

For multi-dimensional inputs, \code{grad o grad} computes the
column of the Hessian corresponding to the second differentiated
argument.

\subsection{grad o vmap}
\label{ssec:higher_order:grad_vmap}

Per-sample gradients:

\begin{lstlisting}[style=lucid,language=LucidPython]
def loss(w, x, y):
    return ((w @ x - y) ** 2).sum()

per_sample_grads = vmap(grad(loss), in_axes=(None, 0, 0))
gs = per_sample_grads(w, x_batch, y_batch)
# gs.shape == (batch, *w.shape)
\end{lstlisting}

The result is a per-sample gradient, useful for adversarial
training, privacy analyses, and certain meta-learning algorithms.

\subsection{vmap o grad}
\label{ssec:higher_order:vmap_grad}

Order matters. \code{vmap(grad(f))} computes the gradient for one
batch element under vmap (i.e., per-sample gradients);
\code{grad(vmap(f))} computes the gradient of the batched
function (the gradient of the sum across the batch, equivalent
to the per-sample gradients summed).

For loss aggregation, \code{grad(vmap(f))} is what you usually
want.

\section{Implementation Notes}
\label{sec:higher_order:impl}

The functional transforms live in \code{lucid/func/}. The
implementations are thin wrappers around the autograd engine and,
where applicable, \MLX's own transforms.

\subsection{Autograd-graph rebuilding}
\label{ssec:higher_order:rebuild}

Higher-order differentiation requires the backward graph itself
to be differentiable. The mechanism is
\code{backward(create\_graph=True)}, which causes the engine's
\code{apply} calls to themselves run under \code{enable\_grad},
building a second-order autograd graph that can be backwarded
again.

The cost of \code{create\_graph=True} is roughly 2$\times$ the
cost of a normal backward, because every backward op is itself
recorded for further differentiation.

\subsection{Per-transform overhead}
\label{ssec:higher_order:per_transform}

Each transform layer adds Python wrapping overhead in the
microseconds-per-call range. For deeply nested compositions
(\code{hessian} is two layers), the wrapping adds up to tens of
microseconds, which is negligible relative to the kernel cost
for any non-trivial input.

\subsection{Limitations}
\label{ssec:higher_order:limitations}

A few:

\begin{itemize}
  \item The transforms work on tensor inputs and outputs. They do
    not handle Python control flow (\code{if/else}, \code{for}
    loops) that branches based on tensor values; the branch
    taken in the forward must be the same in the backward.
  \item \code{vmap} does not work with operations that have data-
    dependent output shape (\code{nonzero}, \code{unique}).
  \item Higher-order differentiation through custom Functions is
    supported but requires the Function's backward to itself be
    differentiable (i.e., to use \Lucid\ primitives, not
    arbitrary Python code).
\end{itemize}

The chapters that follow \chref{ch:gradcheck} document the
behaviour at the limits.

\section{Worked Example: Per-Sample Gradients for MAML}
\label{sec:higher_order:maml}

Model-Agnostic Meta-Learning (MAML, \cite{finn_maml_2017})
requires per-sample gradients through an inner-loop optimisation. The functional transforms
make this expressible naturally.

\begin{lstlisting}[style=lucid,language=LucidPython]
import lucid
from lucid.func import grad, vmap

def task_loss(params, x, y):
    pred = model_fn(params, x)
    return ((pred - y) ** 2).sum()

def inner_step(params, x, y, lr):
    g = grad(task_loss)(params, x, y)
    return _tree_sub(params, _tree_mul(g, lr))

def meta_loss(params, x_train, y_train, x_test, y_test, lr):
    # Inner adaptation
    adapted = inner_step(params, x_train, y_train, lr)
    # Outer evaluation
    return task_loss(adapted, x_test, y_test)

# Per-task meta-gradient via vmap
per_task_meta_grad = vmap(grad(meta_loss),
                          in_axes=(None, 0, 0, 0, 0, None))
meta_grads = per_task_meta_grad(
    params, x_train_batch, y_train_batch,
    x_test_batch, y_test_batch, lr)
# Aggregate
meta_grad = _tree_mean(meta_grads, axis=0)
# Outer optimiser step
params = _tree_sub(params, _tree_mul(meta_grad, meta_lr))
\end{lstlisting}

The composition \code{vmap(grad(...))} computes the per-task
gradient of the post-adaptation loss with respect to the initial
parameters. Each task in the batch sees its own inner adaptation
and its own outer gradient; \code{vmap} parallelises across
tasks; \code{grad} of \code{meta\_loss} threads through the
inner-loop update.

The first-time reader of MAML's algorithm finds it daunting;
the functional formulation is short.

\section{Worked Example: Hessian-Vector Product for CG-Newton}
\label{sec:higher_order:hvp_cg}

A second-order optimiser like CG-Newton needs $H^{-1} g$, the
Newton step. The full Hessian is too large to invert; instead one
uses Hessian-vector products $H v$ inside a conjugate-gradient
iteration to approximately solve $H d = g$ for $d$.

The HVP can be computed in two ways:

\subsection{HVP via grad-of-grad}
\label{ssec:higher_order:hvp_double_grad}

\begin{lstlisting}[style=lucid,language=LucidPython]
def hvp_via_grad(f, x, v):
    # First compute gradient
    def grad_fn(x_):
        return (grad(f)(x_) * v).sum()
    # Differentiate again
    return grad(grad_fn)(x)
\end{lstlisting}

The trick: $\nabla_x (v^\top \nabla_x f) = H v$. This requires
two backward passes (one for the inner gradient, one for the
outer).

\subsection{HVP via JVP-of-grad}
\label{ssec:higher_order:hvp_jvp_grad}

\begin{lstlisting}[style=lucid,language=LucidPython]
from lucid.func import jvp

def hvp_via_jvp(f, x, v):
    def grad_fn(x_):
        return grad(f)(x_)
    _, hvp = jvp(grad_fn, (x,), (v,))
    return hvp
\end{lstlisting}

This computes the directional derivative of the gradient in the
direction of $v$. Mathematically identical to the grad-of-grad
version, but computationally cheaper because forward mode is
cheaper for low-rank operations.

For very large parameter counts, the JVP-of-grad approach is
preferred. \Lucid\ users can pick based on profiling.

\section{Performance Considerations}
\label{sec:higher_order:perf}

The transforms have some overhead beyond a plain backward. We
quantify here.

\subsection{Per-call overhead}
\label{ssec:higher_order:per_call_overhead}

\tabref{tab:higher_order:overhead} reports the framework overhead
per transform layer on M3~Max.

\begin{table}[t]
\centering
\small
\begin{tabular}{lr}
\toprule
Transform & Per-call overhead \\
\midrule
\code{grad(f)}        & ${\sim}10\,\mu$s \\
\code{vjp(f)}         & ${\sim}12\,\mu$s \\
\code{jvp(f)}         & ${\sim}10\,\mu$s \\
\code{vmap(f)}        & ${\sim}25\,\mu$s \\
\code{jacrev(f)}      & ${\sim}M \times$ (single VJP cost) \\
\code{jacfwd(f)}      & ${\sim}N \times$ (single JVP cost) \\
\code{hessian(f)}     & ${\sim}M N \times$ (single VJP) \\
\bottomrule
\end{tabular}
\caption[Functional transform overhead.]{Per-call framework
overhead for each transform on M3~Max. The \code{jacrev},
\code{jacfwd}, and \code{hessian} costs scale with the input
($N$) and output ($M$) dimensions because they internally run
multiple VJP/JVP calls.}
\label{tab:higher_order:overhead}
\end{table}

For typical ML applications where the kernel cost is many
microseconds per primitive, the transform overhead is
negligible. For very small kernels (single-element scalar
gradients), the overhead can dominate.

\subsection{Memory of higher-order gradients}
\label{ssec:higher_order:memory_higher}

\code{backward(create\_graph=True)} doubles the memory cost
relative to a plain backward because the backward operations
themselves are recorded into a second-order graph. For deep
models this is a significant cost; the user should be aware and
use \code{create\_graph=True} only when truly necessary (e.g.,
explicit second-order optimisation).

\section{Summary and Forward References}
\label{sec:higher_order:summary}

\Lucid's functional transforms (\code{grad}, \code{vjp},
\code{jvp}, \code{vmap}, \code{jacrev}, \code{jacfwd},
\code{hessian}) compose on top of the autograd engine to provide
higher-order differentiation. \code{grad o vmap} gives per-sample
gradients; \code{jacfwd o jacrev} gives the Hessian;
\code{vmap o grad} gives batched gradients of a non-batched
function. Each layer adds modest Python overhead; the kernels are
the autograd engine's.

The next chapter, \chref{ch:gradcheck}, treats finite-difference
verification, the standard tool for validating that a custom
backward is correct.

\begin{lucidnote}
For the user: most ML training uses only first-order
\code{backward()}. The functional transforms come up in research
contexts: meta-learning (per-sample gradients), second-order
optimisers (Hessian-vector products), sensitivity analyses
(Jacobians of mid-network activations). For these use cases,
\code{lucid.func} is the API.
\end{lucidnote}
